\uuid{qBWG}
\titre{Division euclidienne}
\auteur{bodin ; hanani ; mzari}
\organisation{exo7}
\licence{CC BY-NC-SA 4.0}
\datecreate{2020-01-20}
\niveau{L1}
\module{Algèbre}
\chapitre{Polynôme, fraction rationnelle}
\sousChapitre{Fraction rationnelle}
\difficulte{2}
\theme{Fractions rationnelles, division euclidienne}
\exo7qcmid{lille-1.polynomes.027}

\contenu{
\begin{question}
\qclassification{105.04}
\qtags{motcle=division euclidienne}
Soit $\frac{P(X)}{Q(X)}$ une fraction rationnelle. On note $E(X)$ sa partie polynomiale (appelée aussi partie entière).
\begin{answers}
    \good{Si $\deg P < \deg Q$ alors $E(X) = 0$.}

    \good{Si $\deg P \ge \deg Q$ alors $\deg E(X) = \deg P - \deg Q$.}

    \bad{Si $P(X) = X^3+X+2$ et $Q(X) = X^2-1$ alors $E(X) = X+1$.}

    \good{Si $P(X) = X^5+X-2$ et $Q(X) = X^2-1$ alors $E(X) = X^3+X$.}
\end{answers}
\begin{explanations}
La partie entière s'obtient comme le quotient de la division euclidienne de $P$ par $Q$.
\end{explanations}
\end{question}
}
